\chapter{室内环境与健康}
\setcounter{chapter}{4}
\chapterintro{
掌握住宅的卫生学意义和基本卫生要求；掌握室内空气污染的来源和特点；掌握室内空气清洁度的评价指标；了解办公场所的卫生学特点及管理要求。
}

% ----- 4.1 住宅卫生 -----
\section{住宅卫生}

\subsection{住宅的卫生学意义}

\begin{defbox}
\textbf{住宅（Residential Building）}：人们为了充分利用自然环境和人为环境因素中的有利作用和防止其不良影响而创造的生活居住环境。
\end{defbox}

\textbf{住宅的卫生学意义：}
\begin{enumerate}[leftmargin=*]
    \item 住宅是人们生活、学习、工作的最重要场所
    \item 住宅的卫生条件和人类健康息息相关
    \item 不良住宅环境对健康不利；良好住宅环境有利于健康
    \item 住宅卫生状况可影响数代人和众多家庭的健康
    \item 住宅环境对健康的影响具有\keyword{长期性和复杂性}
\end{enumerate}

\subsection{住宅的基本卫生要求}

\begin{enumerate}[leftmargin=*]
    \item \imp{小气候适宜}——温度、湿度、气流、热辐射适宜
    \item \imp{采光照明良好}——充足的自然采光和人工照明
    \item \imp{空气清洁卫生}——通风良好，无有害物质积累
    \item \imp{隔音性能良好}——防止噪声干扰
    \item \imp{卫生设施齐全}——上下水、卫生厕所等
    \item \imp{环境安静整洁}——有利于身心健康
\end{enumerate}

\subsection{居室卫生规模}

\textbf{居室卫生规模}是指根据卫生要求提出的居室面积、容积、净高和进深等应有的规模。

\begin{enumerate}[leftmargin=*]
    \item \imp{居室净高}：室内地板到天花板的高度。面积相同时净高越大容积越大，利于采光通风；过低使人压抑且不利于通风散热
    \item \imp{居室面积（居住面积）}
    \item \imp{居室容积}：每个居住者占有的空间容积 = 净高 $\times$ 面积。室内空气中CO$_2$含量是评价空气清洁度的重要指标
    \item \imp{居室进深}：开设窗户的外墙内表面至对面墙内表面的距离。进深与宽度之比不宜 $>$ 2:1，以3:2为宜。\textbf{室深系数}：进深与地板至窗上缘高度之比，一侧采光 $\leq$ 2--2.5，两侧采光 $\leq$ 4--5
\end{enumerate}

\subsection{住宅朝向与间距}

\begin{itemize}[leftmargin=*]
    \item \imp{住宅朝向}：应使主要房间朝南，辅助房间朝北，满足冬季尽量多日照、夏季避免过多日照并利于自然通风
    \item \imp{住宅间距}：前后排建筑物间应保持的最小间隔距离。\textbf{日照间距}：前后两排南向房屋间，为保证后排房屋在冬至日底层获得$\geq$1h满窗日照保持的最小间隔距离
\end{itemize}

\subsection{隔声}

用隔声材料和隔声结构（实体墙板、密封门窗等）阻挡声能传播，将居室相对封闭起来以减少噪声污染。<｜end▁of▁thinking｜>

\subsection{住宅发展新方向}
\imp{健康住宅}：突出健康要素，满足居住者生理、心理和社会多层次需求。\imp{绿色生态住宅}：消耗最少资源和能源、产生最少废弃物，强调节能、节水、节地和治理污染。

\begin{table}[htbp]
\centering
\caption{居室采光评价指标}
\begin{tabularx}{\textwidth}{lXl}
\hline
\textbf{指标} & \textbf{定义} & \textbf{卫生要求} \\
\hline
采光系数 & 室内工作水平面散射光照度与室外全天空散射光照度的百分比 & — \\
投射角 & 室内工作点与采光口上缘连线和水平线夹角 & $\geq$ 27° \\
开角 & 工作点与对侧遮光物上端连线和工作点与采光口上缘连线的夹角 & $\geq$ 4° \\
窗地面积比值 & 采光口窗玻璃面积与室内地面面积之比 & — \\
\hline
\end{tabularx}
\end{table}

\subsection{室内小气候}

\begin{defbox}
\textbf{小气候（Microclimate）}：生活环境中空气的温度、湿度、气流和热辐射等因素对机体热平衡产生明显影响的综合环境条件。
\end{defbox}

\begin{defbox}
\textbf{有效温度（Effective Temperature, ET）}：人体在不同温度、湿度和风速综合作用下产生的热感觉指标。以风速0 m/s、相对湿度100\%、气温17.7°C时产生的温热感为评价标准。

\textbf{校正有效温度（CET）}：将干球温度替换为黑球温度得出的有效温度。

\textbf{湿球-黑球温度指数（WBGT）}：综合反映气温、气湿、气流和热辐射四种物理因素对机体的作用。

\textbf{风冷指数（WCI）}：综合反映寒冷气候下气温和风速对人体温热感的影响，即在低温环境中风速增加产生的冷效应相当于气温下降的度数。

\textbf{热强度指数（TEI）}：根据机体与环境间热交换情况制订的评价指标。
\end{defbox}

\textbf{小气候对人体影响的常用生理指标：}
\begin{itemize}[leftmargin=*]
    \item \imp{皮肤温度（皮温）}：最常用，与温热感觉、脉搏变化基本平行。\textbf{加权平均皮肤温度（WMST）}：测定3--8个点的皮温
    \item \imp{体温}：判断机体热平衡是否受到破坏最直接的指标
    \item 脉搏、出汗量、温热感、热平衡测定
\end{itemize}

% ----- 4.2 室内空气污染 -----
\section{室内空气污染}

\subsection{室内空气污染的来源}

\begin{keybox}
\textbf{室外来源：}
\begin{enumerate}[leftmargin=*]
    \item 室外空气（大气污染物渗入）
    \item 住宅建筑物材料（氡、甲醛等）
    \item 人为带入室内
    \item 相邻住宅污染
    \item 生活用水污染
\end{enumerate}

\textbf{室内来源：}
\begin{enumerate}[leftmargin=*]
    \item \imp{室内燃烧或加热}——燃煤、燃气、烹调油烟
    \item \imp{室内活动}——吸烟、人体代谢产物
    \item \imp{室内建筑装修材料}——甲醛、VOCs、氡
    \item \imp{室内生物性污染}——霉菌、尘螨、细菌
    \item \imp{家用电器}——电磁辐射、臭氧
\end{enumerate}
\end{keybox}

\subsection{室内空气污染的特点}

\begin{enumerate}[leftmargin=*]
    \item 室外污染物对室内空气的污染
    \item 室内外共存同类污染物的叠加暴露
    \item \imp{吸烟}对室内空气污染（最重要的室内污染源之一）
    \item \imp{建筑材料}对室内空气污染（甲醛、氡等）
    \item 空调引起的室内空气污染（"Sick Building"）
\end{enumerate}

\subsection{居室空气清洁度的评价指标}

\begin{table}[htbp]
\centering
\caption{居室空气清洁度评价指标}
\begin{tabularx}{\textwidth}{lX}
\hline
\textbf{指标} & \textbf{意义} \\
\hline
CO$_2$ & 反映室内人员密度和通风状况，$<$0.07\%（清洁） \\
微生物和悬浮颗粒 & 反映室内卫生状况和通风效果 \\
CO & 反映燃烧不完全程度 \\
SO$_2$ & 反映燃煤污染 \\
空气离子 & 负离子浓度反映空气新鲜程度 \\
甲醛 & 反映装修材料污染 \\
\hline
\end{tabularx}
\end{table}

\subsection{保持居室空气清洁度的卫生措施}

\begin{enumerate}[leftmargin=*]
    \item 住宅地段选择（远离污染源）
    \item 建筑材料和装饰材料选择（绿色环保材料）
    \item 合理的住宅平面配置
    \item 合理的住宅卫生规模
    \item 采用改善空气质量的措施（通风、净化）
    \item 改进个人卫生习惯
    \item 合理使用和保养各种设施
    \item 加强卫生宣传教育
\end{enumerate}

% ----- 4.3 病态建筑物综合征 -----
\section{病态建筑物综合征}

\begin{defbox}
\textbf{病态建筑物综合征（Sick Building Syndrome, SBS）}：现代住宅中多种环境因素（如物理、化学因素）联合作用对健康的影响所引起的一种综合征。主要表现为眼、鼻、喉刺激症状，头痛，疲劳，注意力不集中等。与空调系统使用、通风不良、装修材料污染等有关，确切原因尚不清楚。

\textbf{建筑物相关疾病（Building Related Illness, BRI）}：由于人体暴露于建筑物内有害因素（如细菌、真菌、尘螨、氡、一氧化碳、甲醛等）引起的疾病。与SBS的区别在于BRI有明确的病因和典型的临床症状。

\textbf{化学物质过敏症（Multiple Chemical Sensitivity, MCS）}：由于多种化学物质作用于人体多种器官系统，引起多种症状的疾病。特点：复发性、症状呈慢性过程、由低浓度化学污染物质引发、很难找到具体单一对应致病原。
\end{defbox}

% ----- 4.4 办公场所卫生 -----
\section{办公场所卫生}

\begin{defbox}
\textbf{办公场所（Workplace）}：管理或专业技术人员处理某种特定事务的室内工作环境。
\end{defbox}

\textbf{卫生学特点：}
\begin{enumerate}[leftmargin=*]
    \item 办公人员相对集中，流动性小
    \item 办公人员滞留时间长，活动范围小
    \item 分布范围广泛，基本条件和卫生状况相差较大
    \item 存在诸多影响人体健康的不利因素
\end{enumerate}

\textbf{基本卫生要求：}用地选择合理、采光照明良好、小气候适宜、空气质量良好、环境宽松。

% ----- 4.5 复习题 -----
\section{复习题}

\begin{enumerate}[leftmargin=*, itemsep=8pt]
    \item \textbf{【名词解释】}小气候；有效温度（ET）；病态建筑物综合征（SBS）
    \item \textbf{【名词解释】}采光系数；窗地面积比值
    \item \textbf{【简答题】}简述住宅的基本卫生要求。
    \item \textbf{【简答题】}简述室内空气污染的来源及特点。
    \item \textbf{【简答题】}居室空气清洁度的评价指标有哪些？
    \item \textbf{【简答题】}简述保持居室空气清洁度的卫生措施。
    \item \textbf{【简答题】}简述办公场所的卫生学特点。
\end{enumerate}

\subsection*{参考答案要点}

\begin{enumerate}[leftmargin=*, itemsep=5pt]
    \item \textbf{小气候}：生活环境中空气温度、湿度、气流和热辐射等因素的综合环境条件。\textbf{ET}：人体在不同温度、湿度和风速综合作用下的热感觉指标。\textbf{SBS}：现代住宅中多种环境因素综合作用引起的综合征，表现为眼鼻喉刺激、头痛、疲劳等。

    \item \textbf{采光系数}：室内工作水平面散射光照度与同时室外全天空散射光照度的百分比。\textbf{窗地面积比值}：采光口窗玻璃面积与室内地面面积之比。

    \item 六项基本要求：小气候适宜、采光照明良好、空气清洁卫生、隔音性能良好、卫生设施齐全、环境安静整洁。

    \item 来源分室外（大气渗入、建筑材料、人为带入等）和室内（燃烧加热、室内活动、装修材料、生物性污染、家电）。特点：室外污染渗入、叠加暴露、吸烟污染、建材污染、空调污染。

    \item 主要指标：CO$_2$、微生物和悬浮颗粒、CO、SO$_2$、空气离子、甲醛等。

    \item 八项措施：地段选择、材料选择、平面配置、卫生规模、通风净化、个人习惯、设施保养、卫生宣教。

    \item 四项特点：人员集中流动性小、滞留时间长活动范围小、分布广条件差异大、存在多种不利因素。
\end{enumerate}

\newpage

% =====================================================================
%  CHAPTER 5-10 will be appended next...
% =====================================================================

% =====================================================================
%  CHAPTER 5: 水体卫生 【重点章节】
% =====================================================================
\newpage